\label{sec:examples-of-peano-systems}

\begin{topicbox}{Peano Systems as Abstract Structures}
\begin{exposition}
The examples in this section separate the abstract idea of a Peano system from
any particular representation of the natural numbers. A Peano system is not a
specific set such as $\mathbb{N}$. It is a structure consisting of a carrier
set, a distinguished initial element, and a successor operation satisfying the
Peano conditions. Different carriers can therefore realize the same abstract
successor structure.
\end{exposition}

\begin{remark*}[Formal reference]
The formal definition used throughout this section is
\hyperref[def:peano-system]{Peano System}. In each example, the data are
displayed as a triple $(P,S,e)$ consisting of a carrier $P$, a successor map
$S:P\to P$, and a distinguished first element $e\in P$.
\end{remark*}

\begin{remark*}[Interpretation]
The axioms say that the system has one starting point, the successor operation
never merges two different elements, and induction reaches every element of the
carrier. Thus the carrier is an infinite one-way chain generated from its
initial element.
\end{remark*}
\end{topicbox}

